% Chinese paper draft. Development experiment reports live outside latex/.
% !TeX program = xelatex
% Compile twice with XeLaTeX, or use Tectonic. The fallback is TeX Live's Arphic TTF.
\documentclass[11pt]{article}
\usepackage[a4paper,margin=0.82in]{geometry}
\usepackage{fontspec}
\usepackage{xeCJK}
% Use the bundled default Latin fonts; cover CJK roman, sans and mono separately.
\IfFontExistsTF{Noto Sans CJK SC}
 {\setCJKmainfont[AutoFakeBold=1.4,AutoFakeSlant=0.15]{Noto Sans CJK SC}
  \setCJKsansfont[AutoFakeBold=1.4,AutoFakeSlant=0.15]{Noto Sans CJK SC}
  \setCJKmonofont[AutoFakeBold=1.4]{Noto Sans CJK SC}}
 {% TrueType embeds ToUnicode without relying on an external Adobe-GB1 CMap.
  \setCJKmainfont[AutoFakeBold=1.4,AutoFakeSlant=0.15]{gbsn00lp.ttf}
  \setCJKsansfont[AutoFakeBold=1.4,AutoFakeSlant=0.15]{gbsn00lp.ttf}
  \setCJKmonofont[AutoFakeBold=1.4]{gbsn00lp.ttf}}
\usepackage{amsmath,amssymb,amsthm,mathtools}
\usepackage{booktabs,tabularx,longtable,array,enumitem,xcolor,hyperref,fancyhdr}
\usepackage{tikz}
\usetikzlibrary{positioning,arrows.meta}
\definecolor{navy}{RGB}{27,56,81}
\definecolor{lightblue}{RGB}{229,239,248}
\hypersetup{unicode=true,colorlinks=true,linkcolor=navy,citecolor=navy,urlcolor=navy,
 pdftitle={Intervention-Calibrated Distribution Transport for Cross-Embodiment Sim-to-Real Learning}}
\setlength{\parindent}{0pt}
\setlength{\parskip}{0.42em}
\setlength{\emergencystretch}{3em}
\renewcommand{\baselinestretch}{1.10}
\renewcommand{\arraystretch}{1.18}
\newcolumntype{Y}{>{\raggedright\arraybackslash}X}
\newcolumntype{L}[1]{>{\raggedright\arraybackslash}p{#1}}
\setlist{nosep,leftmargin=1.7em,topsep=0.25em}
\renewcommand{\contentsname}{目录}
\renewcommand{\abstractname}{摘要}
\renewcommand{\refname}{参考文献}
\renewcommand{\figurename}{图}
\renewcommand{\tablename}{表}
\pagestyle{fancy}\fancyhf{}
\fancyhead[L]{\small 任务相关行为分布迁移}
\fancyhead[R]{\small 论文草稿}
\fancyfoot[C]{\thepage}
\setlength{\headheight}{18pt}
\newtheorem{proposition}{命题}[section]
\newtheorem{theorem}[proposition]{定理}
\newtheorem{assumption}[proposition]{假设}
\newtheorem{remark}[proposition]{说明}
\renewcommand{\proofname}{推导}
\newcommand{\E}{\mathbb E}
\newcommand{\R}{\mathbb R}
\newcommand{\Prob}{\mathbb P}
\newcommand{\ind}{\mathbf 1}
\newcommand{\cD}{\mathcal D}
\newcommand{\cQ}{\mathcal Q}
\newcommand{\cZ}{\mathcal Z}
\newcommand{\cI}{\mathcal I}
\newcommand{\cF}{\mathcal F}
\newcommand{\cP}{\mathcal P}
\newcommand{\cT}{\mathcal T}
\newcommand{\TV}{\operatorname{TV}}
\newcommand{\tr}{\operatorname{tr}}
\newcommand{\Lip}{\operatorname{Lip}}
\newcommand{\norm}[1]{\left\lVert #1\right\rVert}
\title{\vspace{-1em}\textbf{Intervention-Calibrated Distribution Transport\\for Cross-Embodiment Sim-to-Real Learning}\\[0.6em]
\large 干预校准的任务相关行为分布迁移\\[0.4em]
\normalsize 中文论文草稿}
\author{Draft}
\date{2026 年 10 月 17 日}
\begin{document}
\maketitle
\vspace{-1em}
\begin{abstract}
不同机器人可以共享任务语义，但同一个技能分布在不同执行机制下会产生不同的接触、运动和失败分布。本文研究如何将源任务条件行为分布迁移为目标机器人可执行的行为分布。我们以已有源行为先验为起点，通过有限干预校准任务相关的执行响应，再学习一个 source-conditioned deterministic residual transport。训练始终保留原始任务损失，并只在 matched task-only 近优、公开校准支持内、ensemble disagreement 较低且执行阶段允许的样本上施加 source-response anchoring；独立模型验证上的 task-only fallback 防止明显的模型代理退化。该设计给出的是条件性的 selective anchoring 机制，soft task band 与 fallback 都不构成真实任务或安全保证。理论区分 gate 对训练目标的代数作用、二选一 checkpoint 的模型代理界、事件裕度条件与实际任务表现。实验拟以全 3D、6/7 自由度机械臂为起点，依次研究持件进入受限工位、严格导面接触定位和带恢复的 keyed connector 插入；第四项任务保持插入语义并改变移动身体与支撑耦合。四任务使用九行强方法主表、六格对齐消融、两个机制组和独立训练/校准实例。实机平台证据尚未完成，本文不预先声称部署、普遍性能优势或安全保证。
\end{abstract}
\tableofcontents
\newpage

\section{Introduction}
\label{sec:intro}
不同机械臂可以完成同一定位或插入任务，却具有不同运动学、控制延迟、柔顺性和接触响应。一个源策略描述的不只是单条轨迹，而是对初态、扰动和多种动作选择的条件行为分布。将它部署到另一身体或真实系统后，实际产生的运动、接触和失败概率可能改变。因此，跨 embodiment 与 sim-to-real 可统一看作受目标执行机制约束的任务相关分布迁移问题。

共同状态/动作表示帮助不同身体共享技能\cite{lsa,lad,xskill}；系统辨识与主动探索利用少量真实实验改善部署\cite{asid,spi,optimalexploration}；分布对齐方法已使用联合表示对齐与 OT/UOT 组织 sim-and-real 学习\cite{otreal}。Human2Sim2Robot 进一步展示了物体层任务语义在 embodiment gap 中的作用\cite{human2sim}。这些是本文的方法基础。

表示或指令分布接近，不自动意味着真实交互响应接近；辨识参数准确，也不自动决定源多模态行为应该如何在目标上改变。本文聚焦有限真实干预下，如何构造\emph{由目标机制实现、保留有效任务结构}的行为分布。目标成功轨迹分布通常没有被完整观测，因此它是由任务目标、源行为先验和校准目标机制共同决定的优化对象。

两种机械臂可能具有相近的自由空间轨迹，却在孔口侧向偏移后产生恢复与卡滞两种不同分布。源先验中的多个绕行路径或接触调整方向，也未必全部适用于目标工作空间。迁移需要改变目标行为的实现和概率结构，并允许必要的任务相关偏离；无条件追求两个分布越接近越好并不合理。

从已有源先验采样短时技能，通过少量真实探针校准目标响应，以任务阶段相关的几何定义哪些差异重要，再学习 source-conditioned transport。确定性映射已能保留随机源 prior 的多峰性；只有需要表达额外条件不确定性时才引入 transport noise，并单独验证其收益。训练在校准模型中进行，部署根据当前实际状态直接生成目标技能并闭环执行。action decoding 使用现有机器人接口。

\subsection{Three Research Questions}
\begin{enumerate}[label=\textbf{RQ\arabic*.}]
\item \textbf{Transfer effectiveness.} 在相同目标交互预算下，迁移后的行为分布能否改善跨机械臂与 sim-to-real 的闭环任务表现，并将同一操作技能扩展到移动身体？
\item \textbf{Calibration efficiency.} 有限任务相关干预能否校准影响迁移的执行响应，并在已声明机制族中泛化到未见初态、控制器和接触条件？
\item \textbf{Transport mechanism.} 干预响应、任务几何与可行性约束分别如何影响目标分布的任务对应、多模态覆盖和执行稳定性？
\end{enumerate}

\subsection{Three Candidate Contributions}
\begin{enumerate}[label=\textbf{C\arabic*.}]
\item \textbf{可执行分布迁移的建模。} 将任务条件源行为先验与目标闭环行为分布分开，通过目标执行机制定义 transport 的可实现集合；目标 marginal 是待构造的行为 law。
\item \textbf{干预校准的 realization-aware transport 方法。} 结合任务调制响应几何、低维机制校准、可执行约束和 source-conditioned mapping，学习目标行为分布；映射是否需要额外随机性由消融决定，action decoding 为执行接口，已有生成模型可直接作为 prior。
\item \textbf{预算匹配的分布与任务验证。} 在跨机械臂、接触操作与真实平台上，同时检验闭环效果、干预预算、多模态覆盖和响应分布变化，并用辅助分析解释适用条件。
\end{enumerate}
主要贡献由方法与实验建立。任务相关探索、OT/UOT、flow matching、Lagrange 更新与动作映射均有前例，不单独宣称为新概念。

\section{Related Work}
\label{sec:related}
\subsection{Cross-Embodiment Representations and Behavior Transfer}
Latent Space Alignment 使用编码/解码、cycle consistency 与 latent dynamics 对齐跨机器人技能\cite{lsa}；Latent Action Diffusion 以语义对齐 latent action 支持不同末端执行器的 diffusion policy\cite{lad}；XSkill 研究跨 embodiment 技能发现\cite{xskill}。Human2Sim2Robot 使用物体轨迹构造共同任务奖励\cite{human2sim}。本文继承共同任务语义，关注给定表示后，行为分布在不同机制下的闭环实现。

\subsection{Sim-to-Real Adaptation and Active System Identification}
Dynamics randomization、GAT 和 RMA 分别提供随机化、动作 grounding 与历史适应路线\cite{dr,gat,rma}。ASID 与 SPI-Active 主动采集信息并辨识目标机制\cite{asid,spi}；Guanya Shi 参与的 optimal exploration 已研究任务相关实验设计\cite{optimalexploration}。本文把有限实验得到的机制信息用于构造目标行为 law，而不将 task-aware exploration 本身作为主要创新。

近期 ASAP 在目标数据上学习 delta-action model，并在对齐后的模拟器中
微调源策略\cite{asap}；FADA 将 motion planner 与 inverse dynamics 分离，
冻结前者并用少量目标记录微调后者\cite{fada}。两者原生研究人形动力学适应，
移植到机械臂接触时必须保留完整训练链并明确适配范围。
Play2Perfect 通过多对象 play 预训练和装配微调学习精密接触\cite{play2perfect}，
ACP 从带力示范学习时空可变柔顺控制\cite{acp}。这些方法构成校准后直接学习
任务策略与接触控制的重要竞争路线，不能用简单动作缩放代替完整对照。

\subsection{Distribution Alignment and Transport for Robot Learning}
OT-sim2real 通过联合表示分布对齐、UOT 与时间相关采样进行 sim-and-real co-training，其主要协议使用少量目标演示\cite{otreal}。DARA 与 H2O+ 研究动力学差异下的源/目标学习\cite{dara,h2oplus}。Flow matching 的路径成本和约束结合也已有研究\cite{densityflow}。本文重点在有限干预校准的接触响应、目标可执行性与源行为结构之间的联合建模；不能把联合分布对齐、partial transport 或受约束 flow 单独列为新贡献。

\section{Preliminary}
\label{sec:prelim}
\subsection{Embodiment-Conditioned Control and Action Decoding}
环境 $e$ 有自己的原始观测 $o_e$、动作 $a_e$ 和低层闭环控制。共同状态 $q=\phi_e(h_e)$ 包括物体/TCP 相对位姿、任务阶段、接触与力觉历史；必要时加入机器人能力和配置摘要。不同身体的关节维度不强行混合。

共同技能指令 $U\in\R^{H\times d}$ 表示短时任务空间命令、力设定值及允许的顺应参数。通过已有 action decoding 接口 $\mathcal A_e$ 与反馈控制器 $C_e$ 执行：
\begin{equation}
 a_{t+j}=C_e\big(h_{t+j};\mathcal A_e(U_j,q,c_e)\big),
 \qquad j=0,\ldots,H-1.
 \label{eq:decode}
\end{equation}
编码坐标系、单位、控制时长和允许动作边界需统一。实际接触力为响应，force setpoint 为可设置指令。这里不提出新 decoder 架构，也不把 decoder 单列为创新模块。理论假设共同状态在声明支持内足以描述宏步交互。

移动身体沿用任务坐标系中的同一技能 $U$；$\mathcal A_e,C_e$ 负责将末端参考落地到机械臂、底盘或腿部的全身控制。目标身体状态另包含基座位姿/速度、支撑接触、关节余量与控制器历史，不能只输入 TCP 位姿就假定跨身体状态充分。共享响应对应只比较任务进展、工件/TCP 与任务接触；身体专属的防倾覆、足端/轮地滑移及执行边界进入目标可行性与风险，不要求移动身体复制固定臂的基座运动。已有 UMI-on-Legs 采用 task-frame 末端轨迹连接操作策略与全身控制，并展示固定臂策略向四足平台的迁移；本文需在这一接口基础上验证校准与 transport 的增量价值，而非将接口本身作为新贡献\cite{umilegs}。

\subsection{Conditional Behavior Distributions and Transport}
源 prior 是已有源策略或生成模型：
\begin{equation}
 U_s\sim p_s(U\mid q,g).
\end{equation}
学习 transport kernel $T_\varphi(dU_t\mid U_s,q,g,c_t)$，诱导目标指令分布
\begin{equation}
 p_{t,\varphi}(dU_t\mid q,g)
 =\int T_\varphi(dU_t\mid U_s,q,g,c_t)\,p_s(dU_s\mid q,g).
 \label{eq:kerneltransport}
\end{equation}
确定性特例为 $p_{t,\varphi}=\Phi_\varphi\#p_s$。目标实际行为 law 为
\begin{equation}
 \mathbb P_{t,\varphi}^{g}
 =\operatorname{Rollout}_{K_t,\pi_{t,\varphi}}\#\mu_0.
 \label{eq:pathlaw}
\end{equation}
$p_{t,\varphi}$ 与 $\mathbb P_{t,\varphi}^{g}$ 是不同对象：前者为指令层，后者由目标反馈、动力学和扰动共同实现。源与目标响应使用共同初态/任务条件比较；不把不同初态的边缘分布当作机制差异。

样本重加权可以改变已有支持上的质量；transport 还可以改变目标指令的支持位置。一个可解析例子是 $y_s=u_s$、$y_t=\kappa u_t$：源 prior 位于 $[-1.1,-0.9]\cup[0.9,1.1]$，$\kappa=0.5$ 时，$u_t=2u_s$ 才能保持对应输出。仅在原 compact support 内改权重无法产生该指令分布。若目标边界禁止 $2u_s$，则必须接受不可消除的任务响应偏离。

\subsection{Problem Formulation}
输入包括源行为 prior、共同任务目标、可用目标机器人模型/控制接口，以及总预算 $B$ 的真实校准交互。输出为满足执行约束的目标条件行为分布。目标成功轨迹 law 不预先给定。

给定当前 $q,g$，以固定长度 $H$ 的条件轨迹 law 建立短时 formulation：
\begin{equation}
 \begin{aligned}
 \min_{Q,\Gamma}\quad&
 \E_{\Gamma}\big[C_g(\tau_t)+\lambda c_g(\tau_s,\tau_t)\big],\\
 \text{s.t.}\quad&\Gamma_s=\mathbb P_s^{g,H}(\cdot\mid q),\quad \Gamma_t=Q,
 \quad Q\in\mathfrak P_t^H(q),\\
 &\text{目标执行与风险约束},
 \end{aligned}
 \label{eq:ideal}
\end{equation}
$\mathfrak P_t^H(q)$ 为目标机制及允许短时闭环执行诱导的分布集合。实现中用 $T_\varphi$ 参数化其中的子集，并由源条件与随机种子构造 coupling。部署在每个目标实际状态重新条件化，诱导 Eq.~\eqref{eq:pathlaw} 的全任务 law；拼接的源参考一般不保持独立源 rollout 的全任务 marginal。本文不声称求得所有 coupling 的全局 OT 最优解。

受控协议 A 仅使用声明校准实验，不使用目标任务演示或目标 RL fine-tuning。
补充协议 B-log 允许无 reward/专家示范的计费完整自采轨迹，
B-online 允许计费目标任务交互与 reward，C-supervised 单列目标演示与外部监督，
D-oracle 仅作真值诊断；需要目标演示的 OT-sim2real、LAD 等在其适用监督协议下比较。
所有校准、动作接口拟合和验证计入目标预算 $B$；最终评估不参与训练。

\section{Method: Intervention-Calibrated Distribution Transport}
\label{sec:method}
\textbf{核心流程。} 固定源 prior，建立任务响应表征；用有限可行干预校准目标机制；在校准模型中学习 transport；部署时按当前实际状态采样并执行目标技能。method 的新意需要来自这四步的配合及实证效果。

\begin{figure}[ht]
\centering
\begin{tikzpicture}[node distance=0.6cm and 0.6cm,>=Stealth,
 box/.style={draw=navy,rounded corners,fill=lightblue,align=center,font=\small,
 text width=3.65cm,minimum height=1.02cm}]
\node[box] (prior) {源条件行为分布\\已有策略 / diffusion / flow prior};
\node[box,right=of prior] (cal) {有限真实干预\\校准目标响应与不确定性};
\node[box,below=of prior] (geom) {任务响应几何\\进展、接触与模式变化};
\node[box,below=of cal] (trans) {学习 stochastic transport\\目标机制、任务与可行约束};
\node[box,below=of trans] (exec) {目标行为分布\\action decoding 与闭环执行};
\draw[->] (prior)--(trans);
\draw[->] (cal)--(trans);
\draw[->] (geom)--(trans);
\draw[->] (trans)--(exec);
\draw[->] (exec.west) -- ++(-0.3,0) |- (prior.south);
\end{tikzpicture}
\caption{方法关系。校准模型用于训练 transport；部署依据实际反馈重新条件化源 prior。回路中的重新采样不是在线轨迹 refinement。}
\end{figure}

\subsection{Task-Modulated Response Representation}
\label{sec:geometry}
\textbf{表征对象。} 每个技能执行产生响应描述 $Y=\Psi_g(\tau)$：任务进度、物体/TCP 几何、接触力统计与时间序列、跟踪误差、接触模式及失败事件。源记录的未来轨迹仅用于离线训练；在线 transport 输入只有当前可用历史。

\textbf{任务几何。} 使用有量纲归一化的距离
\begin{equation}
 d_g^2(y,y';q)=(y-y')^\top G_g(q)(y-y'),\qquad
 G_g(q)=\epsilon I+\sum_j\omega_j(q,g)b_jb_j^\top,\quad\epsilon>0.
 \label{eq:metric}
\end{equation}
$\omega_j$ 由阶段、接触状态、目标区域和约束活跃程度指定。受限接近突出全臂/持件间隙与姿态进展，夹具定位突出压靠、滑移与旋转；插入突出横向接触、推进和卡滞/恢复分支。模式概率差异通过 one-hot/事件成本独立进入，不能用平均力掩盖不同失败概率。$\epsilon I$ 避免忽略将影响未来决策的方向；若只保留部分特征，则理论与主张限制在该 quotient space。

\textbf{轨迹成本。} 对相近任务阶段的轨迹定义
\begin{equation}
 c_g(\tau_s,\tau_t)=
 \sum_{k=0}^{H}d_g^2(\Psi_{g,k}(\tau_s),\Psi_{g,k}(\tau_t);q_k^s)
 +\lambda_m c_{\rm mode}(\tau_s,\tau_t).
 \label{eq:semantic}
\end{equation}
首版固定技能时长；明显速度不同的机制可采用受限 phase alignment，并单独保留完成时间成本。不能用任意 time warping 消除真实延迟和力过冲。共同任务成本 $C_g$ 依据进展、接触保持、失败与完成定义；源语义成本与目标任务成本同时存在，允许目标改变源行为以完成任务。

\textbf{多模态结构。} 在 source prior 中保留多个有效模式 $m$，如左右调整或不同绕行路径。第一版用可解释几何规则标记模式，并在目标可执行的模式集合内测量覆盖。mode 标签帮助 transport 避免把多种行为平均成一个无效指令；它不保证所有源模式在目标上都应保留。

\textbf{实现与消融。} 表征使用 TCP/物体状态与 wrist F/T，先不引入大型视觉模型。比较 uniform metric、仅名义几何、平均响应和包含模式/扰动响应的完整表征。这部分吸收 2GO 的任务调制几何 insight。

\subsection{Budgeted Intervention and Target-Mechanism Calibration}
\label{sec:calibration}
\textbf{机制模型。} 在多机制模拟器中预训练短时响应模型
\begin{equation}
 (q',Y)\sim\widehat K_\psi(\cdot\mid q,U,\delta,c),
 \label{eq:model}
\end{equation}
$c$ 为低维执行 context，例如指令增益、延迟、接触顺应性或摩擦参数。真实交互校准 $c_t$ 及预测不确定性：
\begin{equation}
 \hat c_t=\arg\min_c
 \sum_{(q,U,\delta,y)\in\mathcal B_{\rm fit}}
 -\log\widehat K_\psi(y\mid q,U,\delta,c).
 \label{eq:calibrate}
\end{equation}
已有目标模拟器可先校准参数，再训练可微的短时 surrogate。所有目标真实交互（含进入接触阶段、失败和独立验证）计入 $B$，并同时报告控制步数与交互秒数；公开模拟预训练和 calibrated-twin 优化另列模拟计算预算，所有对照具有相同的可用模型与计算额度。少量真实样本用于校准结构化模型，不从头学习任意 world model。物理参数 context 必须限制在声明的物理支持内；离散 delay 用离散假设比较。自由 latent 的坐标不能继续解释为物理参数真值。

\textbf{干预语义。} $\delta$ 是可随机分配的命令、初始偏移或实验配置。定义规定执行接口下的响应
\begin{equation}
 K_e^\delta(dy\mid q,U)
 =\Prob_e(Y\in dy\mid q,\operatorname{do}(U),
 \operatorname{do}(\Delta=\delta);\mathcal A_e,C_e).
 \label{eq:intervention}
\end{equation}
使用匹配初态、已记录历史和随机指派，减少行为选择偏差。第一版探针为自由空间增益/延迟、接触建立、不同切向速度、横向偏移与受限恢复。力监测使用已有执行阈值，不从少量数据推导未知接触的全局安全性。

\textbf{预算选择。} 设探针为有限候选集 $\Delta_F$，每个候选满足 MDOC 式的已知可行性约束。第一版使用已有 task-dependent design，在模型中评估相关响应不确定性的下降：
\begin{equation}
 \min_{\mu}\ \Phi_g(\mu),\qquad
 \E_\mu c(\delta)\le b,\quad\operatorname{supp}\mu\subseteq\Delta_F,
\end{equation}
\begin{equation}
 \Phi_g(\mu)=
 \begin{cases}
 \tr(B_gF_\mu^\dagger B_g^\top),&\ker F_\mu\subseteq\ker B_g,\\
 +\infty,&\text{其他情况}.
 \end{cases}
 \label{eq:probe}
\end{equation}
这是方法支撑，非首次提出任务相关探索。增强版可用 ensemble 的虚拟测量更新，选择降低 transport 成本不确定性最多的探针；它应单独消融，不能与固定设计 Gaussian 公式混用。

\textbf{校准与验证分开。} 拟合集用于机制/接口参数，独立验证集检查各任务模式、响应分布与不确定性。动作接口调整后重新验证；所有拟合、验证和失败探针均计入 $B$。未覆盖的接触分支作为 unknown 报告。source prior 或 transport 改变访问分布时，重新检查模型支持，避免把校准工装覆盖误当作任意任务覆盖。

\subsection{Realization-Aware Distribution Transport}
\label{sec:transport}
\textbf{参数化。} 冻结实现为 source-conditioned deterministic residual map：
\begin{equation}
 U_t=\Phi_\varphi(U_s,q,g,\hat c_t,m)
 =\operatorname{Bound}\big(U_s+\Delta_\varphi(U_s,q,g,\hat c_t,m)\big).
 \label{eq:residual}
\end{equation}
$m$ 表示源行为模式，$\operatorname{Bound}$ 实施已知命令边界。随机性来自
source prior；当前方法不额外注入 transport noise。真实目标响应由 $K_t$
实现，训练预测由 $\widehat K_\psi$ 实现。约束与损失在最终执行指令上评估，
不能仅检查投影前输出。

\textbf{反馈与表达支持。} 源行为实现包含实际接触阶段与位姿误差；阶段由接触、撤离和对准事件推进，不能只用绝对时钟替代。源 prior 与 transport 可条件于可用历史；由 decoder 提供的位姿反馈对所有方法共享。源停止不意味着目标已到平衡点，因此不得用 $|U_s|$ 将目标所有对应纠偏方向硬置零。采用有界残差及实际 $SE(3)$ 反馈维持终止平衡；力控制期间投影法向位置参考，避免未使用参考积累后在释放时产生突变。反馈 decoder 对所有 baseline 相同，其修正不作为 transport 的独立贡献。

conditional flow 可作为未来扩展：
\begin{equation}
 \frac{dX_\ell}{d\ell}=v_\varphi(X_\ell,\ell;U_s,q,g,\hat c_t,m),
 \quad X_0\sim\mathcal N(0,I),\quad U_t=\operatorname{Bound}(X_1).
 \label{eq:flow}
\end{equation}
这是一个 source-conditioned transport kernel；$\ell$ 是生成/transport 时间，
不是物理时间。它尚未进入冻结主方法或主张；若加入，必须与 residual map
使用相同任务目标、数据和预算比较。

\textbf{实际响应目标。} 当前实现不把所有 source response 都强制写入目标。
对一个训练样本，记校准模型中的原始任务成本为 $C_i(\varphi)$，冻结的
matched task-only 策略在同一批样本上的成本为 $C_i^T$，source-conditioned
响应成本为 $R_i(\varphi)\geq0$。定义加法 task-band hinge
\begin{equation}
 H_i(\varphi)=\big[C_i(\varphi)-C_i^T-\delta\big]_+ .
 \label{eq:additive-band}
\end{equation}
任务损失 $C_i$ 在 band 内外都保留；$H_i$ 只是超出 matched task-only
模型代理容差后的附加惩罚，不替代 $C_i$。

响应 anchoring 只在 detached gate 内生效。令 $d_i^{\rm sup}$ 为相对公开
校准支持的距离，$d_i^{\rm ens}$ 为 response ensemble disagreement，
$M_i^{\rm exec}$ 为当前阶段允许改变命令的执行掩码，则
\begin{equation}
 G_i=M_i^{\rm exec}
 \ind\{C_i(\varphi)\le C_i^T+\delta\}
 \ind\{d_i^{\rm sup}\le q_{.95}^{\rm sup}\}
 \ind\{d_i^{\rm ens}\le q_{.95}^{\rm ens}\}.
 \label{eq:selective-gate}
\end{equation}
两个分位阈值只由 source/known-nominal 校准数据确定；$G_i$ 不参与反向传播。
恢复任务把 recovery mask 并入 $M_i^{\rm exec}$，因此 response loss 比较实际
会被恢复控制器采用的 effective action。冻结的核心目标为
\begin{equation}
 J_{\rm gated}(\varphi)=
 \E_i\!\left[C_i(\varphi)+\beta H_i(\varphi)
       +\lambda G_i R_i(\varphi)\right].
 \label{eq:objective}
\end{equation}
实现中 $\delta=0.05$、$\beta=100$；二者都是模型侧训练和选择常数，
不是实际任务容差或安全裕度。训练后只比较候选 checkpoint 与 matched
task-only：若独立 held-out 模型任务代理超过 task-only 加 $0.05$，则直接
回退到 task-only checkpoint。最终目标数据不参与该选择。

完整六格定义为 frozen-warm，$C$，$C+R$，$C+H$，$C+H+R$ 与
$C+H+G R$。其中 $C,C+R,C+H,C+H+R$ 构成严格的无 gate $2\times2$
比较，$C+H+G R$ 与 $C+H+R$ 的差异才单独识别 selective gate。
历史上以 $H$ 替换 $C$ 的 band-only/Full 目标不属于这一实现。

\textbf{从 coupling 到 distribution。} Eq.~\eqref{eq:objective} 对 $U_s\sim p_s$ 与 $U_t\sim T_\varphi(\cdot\mid U_s)$ 的联合 law 优化，其目标 marginal 是 Eq.~\eqref{eq:kerneltransport}。因此它在改变整个目标行为分布及其支持，而非给源 replay 赋权。需要比较的 realized response law 为
\begin{equation}
 \rho_{t,\varphi}^{\delta}(dy\mid q,g)
 =\int K_t^\delta(dy\mid q,U_t)\,
 T_\varphi(dU_t\mid U_s,q,g,\hat c_t)\,p_s(dU_s\mid q,g).
 \label{eq:responselaw}
\end{equation}
为避免 samplewise coupling 过强，可增加批内 Wasserstein/UOT 响应对齐，但它是已有工具的增强版。基准实现先使用可解释的 coupling 和目标任务成本，再比较 nominal feature UOT 与 realized-response UOT。UOT 中未匹配质量不是消失的实际任务；实际策略仍是归一化行为 law。

\textbf{硬约束与接触风险。} 已知运动学、速度/命令变化率和工作空间边界进入可行域。Eq.~\eqref{eq:additive-band} 是 soft penalty，Eq.~\eqref{eq:selective-gate} 是训练掩码；二者均不构成真实系统的不变集或安全证书。模型相关风险可使用 ensemble：
\begin{equation}
 \mathcal R_j(\varphi)=
 \sup_{M\in\mathcal M_{\mathcal B}}
 \operatorname{CVaR}_{1-\delta_j}
 \left[\max_k h_j(\hat q_k,\hat U_k)\right]\le b_j.
 \label{eq:risk}
\end{equation}
必要时用独立误差估计 tightening。该风险是模型条件下的约束，不是实机安全证书；真实执行仍遵守已有验证范围与监测。

\textbf{可选的受约束扩展。} 若未来实现显式 CVaR 约束，可使用标准 Lagrange 乘子协调目标与风险约束：
\begin{equation}
 \mathcal L(\varphi,\lambda)
 =J(\varphi)+\sum_j\lambda_j(\mathcal R_j(\varphi)-b_j),
 \qquad\lambda_j\ge0,
\end{equation}
\begin{equation}
 \varphi\leftarrow\varphi-\eta_\varphi\nabla_\varphi\mathcal L,
 \qquad
 \lambda_j\leftarrow[\lambda_j+\eta_\lambda(\widehat{\mathcal R}_j-b_j)]_+.
 \label{eq:lagrange}
\end{equation}
hard constraints 通过公共执行接口和在线监测实施，soft semantic costs 通过目标体现。当前冻结的比较实现使用 Eq.~\eqref{eq:objective} 与公共监测，不把尚未进入主实验的 Lagrange/CVaR 扩展写成已完成模块。

\textbf{训练实现。} 首版采用短时可微 surrogate，固定 $p_s$ 和校准模型，对 transport 反向传播。模型先以公开训练机制预训练，再用目标 context 校准；隐藏目标真值不参与训练。若可实例化校准 physics twin，可在其中生成同预算数据或用已有 SAC 学习 residual actor，并保持源语义与风险目标。仅有 latent surrogate 时，换优化器仍受同一模型误差限制。两条路线分别报告；不声称黑箱接触天然可微。

\subsection{Training Algorithm and Amortized Closed-Loop Deployment}
\label{sec:deployment}
\textbf{离线/校准阶段的完整算法。}
\begin{enumerate}
\item 使用已有策略训练/固定源条件 prior；保留多个源行为模式与 source dynamics 信息。
\item 在声明候选集内选择 $B$ 次校准交互，划分拟合与独立验证；拟合目标 context 和 ensemble 不确定性。
\item 检查动作接口与模式支持；固定接口，标记验证不足区域。
\item 采样 $q,g,m,U_s,\delta$，通过 $T_\varphi$ 得到 $U_t$；在源与 calibrated target model 中 rollout。
\item 计算 $C,H,G R$，更新 transport；记录 gate coverage、support distance、ensemble disagreement 与任务代理。
\item 在独立模型验证集上执行 matched task-only $+0.05$ checkpoint fallback；锁定方法、校准阈值与拒绝规则后再进行 held-out 目标评估。
\end{enumerate}

\textbf{在线阶段。} 每个技能边界读取实际 $q$，从 $p_s(\cdot\mid q,g)$ 采样，调用 transport，再通过 action decoding 执行固定短时技能；新的真实反馈用于下一次条件化。低层保持闭环。没有每次在线求解轨迹优化的必要。

若不确定性超过预设验证范围，则调用已有 hold/retreat/recovery 接口或停止。该事件单列为覆盖不足，并在未完成任务时计为失败。fallback 不能提供只对本文开放的额外任务监督。

\section{Supporting Analysis}
\label{sec:theory}
本节给出辨识、响应误差和 transport 成本之间的条件链；这些条件本身需要验证，不能由实验成功倒推出成立。

\subsection{Task-Relevant Identifiability}
在固定局部接触模式中，设
\begin{equation}
 y=A_{\mathcal B}\theta+\xi,\qquad b_g=B_g\theta.
\end{equation}
\begin{proposition}
任务响应 $b_g$ 在无噪声校准中可辨识，当且仅当
\begin{equation}
 \ker A_{\mathcal B}\subseteq\ker B_g.
\end{equation}
\end{proposition}
它说明校准应覆盖影响 transport 的机制方向，不一定需要完整物理参数。核空间不满足时不能用伪逆 trace 虚构零不确定性。固定设计 Gaussian 分析可形成响应置信区域；适应性采样需要相应置信工具。该命题属于标准线性辨识，不单独作为原创贡献。

\subsection{Probe Coverage, Identification and Hybrid Events}
\textbf{不可辨识的边界。} 若两个机制对全部允许探针产生完全相同的观测 law，而其任务所需响应或可成功动作集合不同，任何仅使用这些探针的算法都无法保证在两者上正确适应。特别是只在接触时起作用的未知参数，不能仅凭未接触的观测辨识。自由空间辨识可以外推接触的前提是正确的结构模型与对相关参数的充分激励；未接触本身不意味着所有 gain/delay 参数均不可辨识。

\textbf{有限结构族的充分条件。} 对有限候选机制 $\Theta_F$，令固定探针集合的真实/模型响应 signature 为 $z(\theta),\hat z(\theta)$，且真实机制在候选集中。假设观测噪声范数不超过 $\eta$，所有候选的 probe-model 偏差不超过 $\epsilon_p$。若
\begin{equation}
 \min_{\theta\ne\theta'}\|\hat z(\theta)-\hat z(\theta')\|>2(\eta+\epsilon_p),
\end{equation}
则 nearest-signature 选择恢复真实候选：真实候选到观测的距离不超过 $\eta+\epsilon_p$，其余候选由 triangle inequality 严格更远。连续参数的离散化还需加入覆盖误差；有限搜索不自动获得全局优化或统计置信保证。此结论解释一个受限结构模型的可验证机制，不声称辨识任意接触动力学。

\textbf{从校准到任务的缺口。} Probe 上的小误差没有直接推出下节的 uniform joint-kernel error。必须另声明结构模型正确性，或在 transport 访问支持上验证误差与覆盖；支持外概率与模型偏差不能省略。独立验证阈值是经验拒绝规则，不能称为概率安全证书。

\begin{proposition}[结构族内的校准到响应桥接]
令 $d_\Theta$ 为已声明机制族上的距离（可含离散 delay），真实机制为 $\theta_\star$。假设探针模型 signature 在该族上满足
\[
 \|\hat z(\theta)-\hat z(\theta')\|\ge\kappa d_\Theta(\theta,\theta'),\qquad \kappa>0.
\]
若观测噪声不超过 $\eta$，真实机制处 probe-model 偏差不超过 $\epsilon_p$，所选估计的实际拟合残差为 $r=\|\hat z(\hat\theta)-z_{\rm obs}\|$，则
\[
 d_\Theta(\hat\theta,\theta_\star)\le(r+\eta+\epsilon_p)/\kappa.
\]
进一步假设在共同访问支持及所有允许 transport 上，联合模型 kernel 对机制统一为 $L_\Theta$-Lipschitz，真实机制处的统一 kernel 偏差不超过 $\epsilon_K$，则下一节的误差可取
\[
 \varepsilon_{\mathcal B}\le\epsilon_K+\frac{L_\Theta}{\kappa}(r+\eta+\epsilon_p).
\]
\end{proposition}
\begin{proof}
第一式由 inverse separation 和经观测的 triangle inequality 得到；第二式在真实 kernel 与估计模型 kernel 之间插入真实机制处的模型 kernel，再应用 triangle inequality。该残差形式不要求优化器达到全局最优。
\end{proof}
这补齐了一个\emph{条件性}链条：结构族与任务支持正确、signature 可分、参数到联合闭环响应稳定，才可从探针外推到 transport。这些条件中的 $\kappa,L_\Theta,\epsilon_K$ 需独立论证；有限验证不足以建立其全局界。跨接触事件时统一 Lipschitz 条件可能不成立，应改用局部模式内界与下述事件 margin，不能把连续 gain 的小误差直接称为全任务保证。

\textbf{接触事件。} 设真实/预测轨迹有共同 coupling，轨迹距离的期望不超过 $e$。对一个成功/失败事件，若真实轨迹到事件边界的 margin 为 $M$，则任意 $r>0$ 下
\begin{equation}
 |\Pr(A)-\Pr(\hat A)|\le e/r+\Pr(M\le r).
\end{equation}
这是事件不一致只可能发生于距离至少 $r$ 或 margin 不超过 $r$，再应用 Markov inequality 的结果。孔口与力阈值附近的 margin 项可能很大；连续响应均值误差不能直接替代事件误差。接触建立、卡滞、恢复和超力需分别报告。

\subsection{Calibration Error and the Transport Objective}
分析给定初态的短时 transport 目标。令增广状态 $x$ 包含目标状态/历史、controller carry、事件 phase 与 tracking integral、源参考、源/目标技能指令和模式；真实与预测 rollout 使用同一个源参考及 transport sample。源参考仅是分析与训练成本变量，不向在线控制输入未观测目标未来。假设在同一个共同可行参数集合 $\Phi$ 内，所有 $\varphi$ 的联合增广闭环 kernel 都统一满足
\begin{equation}
 \sup_{\varphi\in\Phi,\,x\in\mathcal X_\Phi}
 W_1(\mathbb K_t^\varphi(x),\widehat{\mathbb K}_t^\varphi(x))
 \le\varepsilon_{\mathcal B},
\end{equation}
输出包含所用任务响应；支持覆盖与联合 kernel 的 $L$-Lipschitz 常数也在 $\Phi$ 内统一成立。仅有目标 marginal 接近不足以控制 source-coupled cost。相同增广初态下分布误差满足
\begin{equation}
 d_{k+1}\le Ld_k+\varepsilon_{\mathcal B},\qquad
 d_k\le\varepsilon_{\mathcal B}\sum_{j=0}^{k-1}L^j.
 \label{eq:propagation}
\end{equation}
对于由期望任务成本和 source-coupled trajectory cost 组成的目标，若相关成本具有统一 Lipschitz 常数，可得到
\begin{equation}
 |F_t(\varphi)-\widehat F_t(\varphi)|\le\varepsilon_F.
\end{equation}
在上述统一条件和共同可行集合中，若模型短时目标训练为 $\varepsilon_{\rm opt}$ 次优，则
\begin{equation}
 F_t(\hat\varphi)-\inf_{\varphi\in\Phi}F_t(\varphi)
 \le2\varepsilon_F+\varepsilon_{\rm opt}.
 \label{eq:objectivebound}
\end{equation}
这是所定义 transport objective 的辅助保证，包含 source anchoring 的偏置；不是对完整目标 task-only optimal policy 的 regret。CVaR、罕见事件和 UOT 增强目标需另外分析，不能直接由此推出真实安全或通用 OT 保证。

\textbf{改进何时可保证。} 对同一可行集合内的基线 $\varphi_0$，若模型成本改进 $\widehat F_t(\varphi_0)-\widehat F_t(\hat\varphi)>2\varepsilon_F$，则真实成本严格改进，直接由两端的统一成本误差界相减得到。若该差值不足，则上述分析不保证优于 direct transfer 或 system identification；随机性、flow 或 OT 名称本身也不产生改进保证。

\subsection{Feasibility Before Response Alignment}
\label{sec:finite-feasibility}
\textbf{校准、表达与选择。} 校准误差、控制/映射表达误差和目标选择偏置
必须分开。小的响应预测误差可以准确预测一个坏策略；
它既不保证候选类存在可行策略，也不保证 source anchoring 优于 task-only。
以下结论条件于固定接口、已声明模型与候选类。

\textbf{可行支持与拒绝。} 设 $\Phi_F$ 是冻结的有限映射族，
$S_m\subseteq\Phi_F$ 是模式 $m$ 下满足所声明任务、约束及模式对应的映射集合。
必须区分真实的 $S_m$ 与仅在有限模型训练/验证 reset 上观察到的
$\widehat S_m$；后者不是真实安全集合的证书。
\begin{proposition}[共同映射、模式条件映射与重加权的支持边界]
对于具有正源概率的模式集合 $\mathcal M_s$，一个共同映射能保留全部模式
当且仅当 $\bigcap_{m\in\mathcal M_s}S_m\ne\varnothing$。
允许逐模式选择同一候选族中的映射时，全部模式可执行当且仅当每个
$S_m\ne\varnothing$。仅对固定映射的源模式重加权不能使其不可执行模式
变为可执行；若全部模式均无支持，则输出停止/拒绝质量，而不是归一化零权重。
\end{proposition}
\begin{proof}
共同映射必须同时属于每个集合。模式条件映射从每个非空集合选择一个元素，
反向必要性直接成立。重加权只改变固定响应支持上的概率，不产生新的命令或
闭环响应支持，故不能补回缺失的可执行模式。
\end{proof}
这是\emph{候选类内}的结论；搜索失败不证明机器人全局不可行。
强 task-only 对照应允许相同逐模式选择，因此增加条件化容量本身不能
作为 response anchoring 的独立优势。有限初态上的全成功过滤仅是经验
拒绝规则；所有拒绝与未执行任务保留在总体分母中。

\begin{proposition}[不变执行前缀的失败上界]
设一族 transport 在到达可观测阶段边界 $\sigma$ 前，采用相同指令、
相同控制器及反馈规则。将初态、源样本和物理随机流共同耦合，则它们在
$\sigma$ 前具有相同的增广闭环轨迹。若共同前缀在到达 $\sigma$ 前
以概率 $p_{\rm pre}$ 触发不可撤销终止失败，则该族任一策略的整体成功率
至多为 $1-p_{\rm pre}$。
\end{proposition}
\begin{proof}
相同初态和随机流下，按闭环递推归纳，相同反馈指令给出相同状态、
controller carry 与事件。故前缀失败事件对全族相同；整体成功事件必须
包含于其补集。对该包含关系取概率即得。
\end{proof}
这里的 $p_{\rm pre}$ 是声明初态分布下的概率，有限模型样本只给出其经验
描述。若不可恢复失败发生在映射开始作用之前，仅优化后续阶段无法消除
该失败。对恢复前缀扩大控制支持时，应向同类控制对照开放相同参数化，
并保持任务约束不变。

\textbf{保留模式质量与停止的代价。} 设 $S$ 是给定可行性规则允许保留
的源模式集合，$p_s(S)>0$。若目标 law $q$ 只在 $S$ 上分配质量，则
\[
 \operatorname{TV}(q,p_s)\ge p_s(S^c),\qquad
 \min_{q:q(S)=1}\operatorname{KL}(q\Vert p_s)=-\log p_s(S).
\]
第二个最小值由 $q=p_s(\cdot\mid S)$ 达到：对任意支持于 $S$ 的 $q$，
$\operatorname{KL}(q\Vert p_s)=
\operatorname{KL}(q\Vert p_s(\cdot\mid S))-\log p_s(S)$。
第一个式子直接在集合 $S^c$ 上使用 TV 定义。若保留原源采样而明确对
$S^c$ 的质量停止，则该停止策略的成功率至多为 $p_s(S)$。
这些是支持限制的基本恒等式，不作为新的 OT 理论；实验中仅有模型估计
$\widehat S$，不能把它当成真实全局可行集合。特别是未通过有限训练规则
不等于该模式成功概率为零；上述停止上界描述算法选择，而不是机器人能力上界。

因此，部分能力条件下不能同时要求源 law 不变、全部模式对应和高任务成功。
重加权应按任务完成、模式请求和分布变化共同评价，不能仅因原始源样本的
对应下降就判为错误。若共同优化模式 law 与映射，应把 $(q,\phi)$ 作为联合候选；
下述固定 law/同集合的任务预算保证不能自动转移到不同重加权策略。

\textbf{用任务非劣预算组织响应目标。} 固定一个非空的共同候选集合
$\mathcal A$，令真实任务与响应成本分别为 $C,R$，模型成本为
$\widehat C,\widehat R$。先求 $\widehat C_*=\min_{\phi\in\mathcal A}
\widehat C(\phi)$，再选择
\begin{equation}
 \widehat\phi_\delta\in\arg\min_{\phi\in\mathcal A}\widehat R(\phi)
 \quad\text{s.t.}\quad
 \widehat C(\phi)\le\widehat C_*+\delta.
 \label{eq:task-band-transport}
\end{equation}
任务约束、mode 条件、候选容量和模拟缓冲对完整方法与 task-only 相同。
$\delta$ 在封存目标评估前固定；它公开允许的任务成本代价，而不是通过
反复调节响应权重寻找最终测试优势。

\begin{proposition}[同一有限集合内的任务预算与响应比较]
若在 $\mathcal A$ 上 $|C-\widehat C|\le\epsilon_C$、
$|R-\widehat R|\le\epsilon_R$ 统一成立，且 Eq.~\eqref{eq:task-band-transport}
求解误差至多为 $\epsilon_{\rm opt}$（约束仍被满足），记
$C_*=\min_{\phi\in\mathcal A}C(\phi)$，则
\begin{equation}
 C(\widehat\phi_\delta)\le C_*+\delta+2\epsilon_C.
\end{equation}
若 $\delta\ge2\epsilon_C$，还有
\begin{equation}
 R(\widehat\phi_\delta)\le
 \min_{\substack{\phi\in\mathcal A\\C(\phi)\le C_*+\delta-2\epsilon_C}}
 R(\phi)+2\epsilon_R+\epsilon_{\rm opt}.
\end{equation}
\end{proposition}
\begin{proof}
以真实 task minimizer 代入模型最小值，得到
$C(\widehat\phi_\delta)\le\widehat C_*+\delta+\epsilon_C
\le C_*+\delta+2\epsilon_C$。任何满足第二式真实成本约束的候选均满足
$\widehat C(\phi)\le C_*+\delta-\epsilon_C
\le\widehat C_*+\delta$，因而属于模型预算带；再对响应成本两端应用统一
误差界和优化误差即得。
\end{proof}
该命题只比较同一候选集合。容差 $\delta$ 须在独立评估前确定；若未认证
统一 $\epsilon_C,\epsilon_R$，经验预算带不能被解释为总体非劣或真实安全保证。

\textbf{冻结实现的 selective gate 只改变 anchoring 项。}
对固定参数 $\varphi$ 和固定 detached gate $G\in\{0,1\}$，记
$H=[\widehat C-\widehat C_T-\delta]_+$。加法 band 的 gated 与 ungated
模型目标分别为
\[
 \widehat J_G(\varphi)=\E[\widehat C+\beta H+\lambda G\widehat R],
 \qquad
 \widehat J_U(\varphi)=\E[\widehat C+\beta H+\lambda\widehat R].
\]
由于 $\widehat R\ge0$，有精确恒等式
\begin{equation}
 \widehat J_G(\varphi)=\widehat J_U(\varphi)
 -\lambda\E[(1-G)\widehat R]\le \widehat J_U(\varphi).
 \label{eq:gated-ungated-identity}
\end{equation}
它只说明同一个 $\varphi$ 下关闭 unsupported anchoring 会降低所定义的训练目标；
两种目标各自优化后的策略不同，Eq.~\eqref{eq:gated-ungated-identity} 不给出
真实任务成本、成功率或安全性的次序。特别地，gate coverage 低可能只是把
方法退化为 task-first 行为；coverage 高也不说明被接受的 source response
对目标任务有帮助。

\begin{proposition}[二选一 checkpoint fallback 的条件任务界]
设 $\widehat C_T,\widehat C_A$ 分别是 matched task-only 与 anchoring
候选在冻结 held-out 模型代理上的成本，选择规则在
$\widehat C_A\le\widehat C_T+\delta$ 时选 $A$，否则选 $T$。
若只对这两个冻结 checkpoint 有
$|C_T-\widehat C_T|\le\epsilon$ 与
$|C_A-\widehat C_A|\le\epsilon$，则被选 checkpoint $S$ 满足
\begin{equation}
 C_S\le C_T+\delta+2\epsilon.
 \label{eq:fallback-pair-bound}
\end{equation}
\end{proposition}
\begin{proof}
若回退到 $T$，结论由 $\delta,\epsilon\ge0$ 直接成立。若选择 $A$，则
$C_A\le\widehat C_A+\epsilon\le\widehat C_T+\delta+\epsilon
\le C_T+\delta+2\epsilon$。
\end{proof}
该结论只需二点误差界，不需要宣称覆盖整个神经网络参数类；但实验仍未认证
真实任务代理误差 $\epsilon$。因此实现中的 $+0.05$ 只能称为模型选择容差，
Eq.~\eqref{eq:fallback-pair-bound} 不能改写成真实 success 非退化或约束保证。
soft hinge 训练期间的瞬时 band violation 也不会被该二选一结论消除。

\textbf{gate 帮助事件对应所需的额外条件。} 对接触切换、卡滞与恢复等
hybrid event，至少需要同时满足三项：被 gate 接受的 anchoring 在真实目标上
具有正确方向；校准模型与数值实现误差小于事件判据的轨迹裕度；该裕度在
增广状态中包含 phase、guard counter、integral、recovery mask 与其他
controller carry。若任一条件不成立，减小平均 response discrepancy 仍可
跨过错误事件边界并降低任务成功。故 selective gate 是一个可检验的
条件机制，不预测跨任务一致的正收益。

\textbf{模型偏差与 reset 采样误差分开。} 上式的 $\epsilon_C$ 若将有限
训练均值当作 $\widehat C$，必须同时覆盖模型偏差与 Monte Carlo 误差。
对事先固定的 $M$ 个候选和校准后独立的 $n$ 个 reset clusters，
若每 cluster 的模式平均成本属于 $[0,1]$，则以概率至少 $1-\alpha$，可以取
\[
 \epsilon_C=\epsilon_{\rm model}
       +\sqrt{\frac{\log(2M/\alpha)}{2n}}.
\]
此处 $\epsilon_{\rm model}$ 仍需外部论证；有限 occupancy 的实测误差
不是它的统一证书。即使 twin 精确预测实际轨迹，有限训练初态上的最优候选
也可能在新的独立初态上失败。当采样误差大于所选任务容差时，经验筛选
不能支持总体非劣结论。
对有限样本的精确枚举最优，也不等于对期望目标的 $\epsilon_{\rm opt}=0$。

\textbf{数值模型误差单独计入。} 设 $C_{\rm ref}$ 为目标参考机制的
期望成本，$C_{\rm twin,ref}$ 为校准模型在参考数值实现下的成本，
$C_{\rm twin,train}$ 为同一模型在训练数值实现下的成本。若分别存在同一有限
候选集合上的统一界，则由三角不等式，模型训练均值所需的统一误差可取
\[
 \epsilon_C = \epsilon_{\rm calibration}
                  +\epsilon_{\rm engine}
                  +\epsilon_{\rm sampling},\qquad
 \epsilon_{\rm engine}\ge
 \sup_{\phi\in\mathcal A}
 |C_{\rm twin,ref}(\phi)-C_{\rm twin,train}(\phi)|.
\]
这里各项是需要论证的上界，而非由少数配对轨迹自动取得的证书。
接触求解的数值差异可能通过阶段切换、保持终止或不可恢复失败放大。
因此校准精度与训练模型的任务预测精度是不同条件；有限配对验证也不能
被解释为对所有候选与未来初态的统一界。

\textbf{表达误差不能被模型精度消除。} 对更宽的参考策略类
$\Psi\supseteq\mathcal A$，定义
$\Delta_{\mathcal A}=\inf_{\phi\in\mathcal A}C(\phi)-
\inf_{\psi\in\Psi}C(\psi)\ge0$。若原标量化目标
$\widehat C+\lambda\widehat R$（$\lambda\ge0$）的优化误差为 $\epsilon_{\rm opt}$，
$R\ge0$，则以同集合真实 task minimizer $\phi_C$ 为比较点有
\begin{equation}
 C(\widehat\phi)-\inf_{\psi\in\Psi}C(\psi)
 \le\Delta_{\mathcal A}+2\epsilon_C+2\lambda\epsilon_R
   +\lambda R(\phi_C)+\epsilon_{\rm opt}.
\end{equation}
证明由模型标量化最优性、成本两端误差以及
$-\lambda R(\widehat\phi)\le0$ 直接得到。即使校准和优化完全准确，
候选表达误差与 anchoring 偏置仍可能存在；上一节的
$2\varepsilon_F+\varepsilon_{\rm opt}$ 并未消除这些项。

\textbf{接触事件裕度与策略质量分开。} 上节事件界控制真实与预测的
事件差异，不保证事件本身成功。应同时报告任务失败概率、逐约束轨迹裕度、
以及预测事件误差。均值误差、最终力或模式标签不足以表示执行鲁棒性。
事件距离必须包含接触状态、计数器与控制器记忆；
末端位置或平均力的 RMSE 不能替代它。
同时，控制器在首次满足保持条件时终止，会使最后保持窗口的起点接近阈值。
因此上述窗口裕度只是轨迹描述，不能直接当作闭环扰动鲁棒半径；验证后者
还需要控制扰动、完整停止时刻与必要的终止后承载诊断。

\textbf{验证拒绝属于算法成本。} 令 $a_{j,m}\in\{0,1\}$ 为目标评估前冻结
的模式接受标记，$C_{j,m}\in[0,1]$ 为该映射执行时的真实任务成本。
以拒绝成本 $1$ 计入全部请求时，实际算法成本是
\[
 C_j^{\rm alg}=\sum_m p_m\{a_{j,m}C_{j,m}+1-a_{j,m}\}.
\]
因此同一源 law 下，完整方法 $F$ 与 task-only $T$ 满足
\[
 C_F^{\rm alg}-C_T^{\rm alg}
 \le \sum_{m:a_{F,m}=a_{T,m}=1}p_m(C_{F,m}-C_{T,m})
       +\sum_{m:a_{F,m}=0,\,a_{T,m}=1}p_m.
\]
共同拒绝模式贡献为零；仅 task-only 被拒绝的项非正；仅完整方法被拒绝
的项为 $p_m(1-C_{T,m})\le p_m$，即得该式。训练预算带即使控制了
映射本身的成本，也不能忽略两方法的验证拒绝差异。若同时采用不同的
模式权重 $p_F,p_T$，按 $p_F$ 应用上述比较后还需加上
$\operatorname{TV}(p_F,p_T)$，因为停止后的成本仍属于 $[0,1]$。
这区分了被执行映射的质量、接受集合和模式质量分配，防止只在成功部署
子集上宣称总体非劣。若两方法具有相同冻结映射、接受标记与源 law，则
其算法成本由结构直接相等；这不提供独立任务收益。

\textbf{边缘一致不等于源条件对应。} 令 $M\sim\mathrm{Bernoulli}(1/2)$。
响应 $Y=M$ 与 $\widetilde Y=1-M$ 的边缘 law 完全相同，但
$\Pr(Y\ne M)=0$、$\Pr(\widetilde Y\ne M)=1$。
两种策略可以具有相同成功率与目标模式频率，却具有相反的联合对应。
这一构造说明评估应区分条件对应与边缘分布；它不能单独证明
本文方法优于使用同一 source prior 的 task-only 控制器。

\textbf{有限类选择持平的解释。} 对冻结成本 $(C_i,R_i)$，候选 $i$
最小化 $C+\lambda R$ 当且仅当
$(C_i-C_j)+\lambda(R_i-R_j)\le0$ 对全部 $j$ 成立。
这些半直线的交集给出可直接计算的 $\lambda$ 区间，包含空区间和无界区间。
同一个映射在共同随机流下产生相同实际 law；给它不同方法名称不会产生
独立收益。该代数诊断说明何时 anchoring 不改变选择，既不是统计证书，也
不支持通过试遍 $\lambda$ 后挑选测试集胜者。

\textbf{相同控制类与额外随机性的边界。} 本文 full 与强 task-only 使用
相同参数化与信息，因此可表达相同闭环 law；source anchoring 改变的是
选择准则，不扩大相对该强对照的可执行集合。相对 reweight-only 的支持
改变不能转称相对同类控制器的表达优势。

固定一个源样本、可观测条件和有限合法动作集合，令 $c(a)$ 为已对物理
随机性取期望的标量成本。任意条件随机 kernel $T$ 满足
$\sum_a T(a)c(a)\ge\min_a c(a)$，确定性最小动作即可达到等号。
因此在可逐条件分解的期望目标中，额外 transport noise 没有自动收益；
随机源 prior 经确定性映射仍可多模态。分布耦合约束、条件覆盖需求或
不同的鲁棒博弈可能改变这一结论，必须单独分析。随机与确定性映射应在
相同容量、目标及训练预算下比较；固定附加扰动不能替代经过训练的随机 kernel。


\subsection{Finite-Policy Evaluation Bounds}
对预先冻结的有限策略集合，可以采用有限样本的成本比较，而不宣称验证了上述无限策略类的统一 kernel 假设。设策略数为 $M$，校准后独立抽取 $n$ 个 reset clusters，每个 cluster 内平均所有预声明模式；同一 cluster 的实际与模型成本使用匹配随机流。成本函数在评估前固定并有界于 $[0,1]$，记观测为 $a_{ji},p_{ji}$，其均值为 $\bar a_i,\bar p_i$。令
\[
 \beta_n=\sqrt{\frac{\log(4M/\alpha)}{2n}},\qquad
 \epsilon_n=\min\!\left\{1,\max_i\frac1n\sum_j|a_{ji}-p_{ji}|+
 \sqrt{\frac{\log(2M/\alpha)}{2n}}\right\}.
\]
\begin{proposition}[有限冻结策略的成本改进界]
若 clusters 独立同分布且与策略选择数据独立，则以至少 $1-\alpha$ 的概率，对所有冻结策略 $i$，
\[
 F_t(\varphi_0)-F_t(\varphi_i)
 \ge\bar p_0-\bar p_i-2\epsilon_n-2\beta_n.
\]
\end{proposition}
\begin{proof}
对 $M$ 个模型均值使用双侧 Hoeffding 界与 union bound，对 $M$ 个有界绝对配对误差使用单侧 Hoeffding 界与 union bound，各分配 $\alpha/2$。由 $|F_t(\varphi_i)-\widehat F_t(\varphi_i)|\le\mathbb E|a_i-p_i|$，再对基线和策略两端应用 triangle inequality 即得。
\end{proof}
该结论条件于一次校准、声明的 reset 分布、成本和有限策略；多次校准需另外分层或分配置信预算。小样本时右端可能为负，必须返回未认证。该界不能覆盖未冻结的 neural policy、支持外接触、任意实机噪声或任务成功指示函数之外未记录的风险，也不保证相对采用不同目标的控制器必胜。

\paragraph{直接观测的成本改进。}
若只验证冻结策略的实际效果而不认证模型，可以对独立 reset 的配对成本差 $d_{ji}=a_{j0}-a_{ji}\in[-1,1]$ 使用更直接的单侧界。对 $M-1$ 个预声明比较，以至少 $1-\alpha$ 的概率同时有
\[
 F_t(\varphi_0)-F_t(\varphi_i)
 \ge \bar d_i-\sqrt{\frac{2\log((M-1)/\alpha)}{n}}.
\]
证明直接对长度为 $2$ 的有界配对差应用 Hoeffding 与 union bound。该式只认证固定成本的实际改进，不能代替校准可辨识性、response kernel 精度或跨校准统计；样本量与比较集合须在这些新 reset 执行前固定，禁止 optional stopping。

\section{Experiments}
\label{sec:experiments}
实验围绕三个问题组织：少量目标交互是否改善受约束任务表现；校准后的
分布迁移是否提供超出动作适应与模型控制的收益；这种收益能否推广到未见
机制和身体。本节给出拟采用的比较协议，正式结果表和结论留待独立评估。

\subsection{Tasks, Bodies, and Shared Control Interface}
源机器人为 Panda，固定臂目标为 xArm7，并以 UR5e 检验跨自由度迁移。
目标几何、关节边界、工具与控制能力对全部方法公开；未知控制响应和接触
参数通过声明的目标交互获得。身体变化与机制变化独立交叉，分别比较同身体
不同响应、不同身体相近响应，以及身体与响应同时变化。

四项任务使用任务空间技能 $U\in\R^{H\times13}$：
\begin{equation}
 U_j=\big[v_j^\top,\ \omega_j^\top,
 (\log k^{\rm trans}_j)^\top,
 (\log k^{\rm rot}_j)^\top,\ F^d_{n,j}\big]^\top.
 \label{eq:experimentalprimitive}
\end{equation}
其中 $v,\omega$ 为线速度与角速度，$k$ 为相对参考值的对角刚度，
$F_n^d$ 为法向力设定值。各身体通过自己的低层控制接口实现这些指令；
自由空间阶段屏蔽力指令。任务阈值、坐标系、控制周期、碰撞白名单、
终止规则和初态分布在开发阶段确定，并在正式评估前冻结。

\begin{enumerate}[label=\textbf{Task \arabic*.},leftmargin=2.1em]
\item \textbf{ObstacleTransfer.} 持带方向性的刚性工件绕过横梁与侧板，
穿过姿态受限的三维开口，到达指定 $SE(3)$ 位姿。源行为包含上绕与侧绕
路径；评估几何差异、执行延迟与负载变化对无碰撞完成、制动和跟踪的影响。
\item \textbf{GuidedContactAlignment.} 沿三维导面建立接触，通过滑移和旋转
到达定位区域。比较控制响应、摩擦与柔顺性变化，重点检验自由空间跟踪
相近时，是否仍需调整压靠、速度、刚度和接触序列。
\item \textbf{KeyedInsertionRecovery.} 从横向与角度偏移启动，将非轴对称
connector 插入匹配 socket；包含受限搜索、回撤、重对齐和再次推进。
重点评价整体完成、卡滞后的恢复及瞬时侧向力，而非仅评价平均接触力。
\item \textbf{MobileKeyedInsertionRecovery.} 保持 Task 3 的工件、夹具、
力限和恢复语义，改变身体支撑与接触耦合。检验固定臂技能在移动机械臂
和带臂腿式平台上的受约束实现。
\end{enumerate}

\begin{table}[ht]
\centering\small
\begin{tabularx}{\linewidth}{L{2.1cm}YY}
\toprule
任务 & 主要任务指标 & 机制与鲁棒性指标\\\midrule
Task 1 & 无碰撞完成、位姿误差、时间 & 全臂/持件间隙、跟踪误差、饱和\\
Task 2 & 限力定位、定位误差、时间 & 接触丢失、滑移/停滞、力峰值\\
Task 3 & 受约束插入、深度、时间 & 卡滞、条件恢复、重试、侧向力峰值\\
Task 4 & 插入且身体约束满足、总时间 & 支撑失败、滑移、身体偏移、WBC 饱和\\
\bottomrule
\end{tabularx}
\caption{四任务的比较重点。各任务均保留全部请求的失败与拒绝分母。}
\end{table}

\subsubsection{Task 4: Same Operation with Mobile Support}
\label{sec:task4}
Task 4 复用冻结的 Panda Task 3 prior。先比较固定与移动的同一机械臂，
再检验带臂腿式身体；真实轮地或足地动力学与平台约束分别定义。
平面运动代理只能用于接口诊断。跨移动支撑类型的结论需要多个有效支撑
类型的独立证据；人形平台取决于实际控制接口，不作为预设完成项。

同一身体上的全部高层方法共享 WBC、站位规则和可行性投影。
任务条件包含可用的基座姿态/速度、支撑接触、关节余量与控制器历史，
身体专属状态不直接塞入冻结的 Panda prior。末端响应对应允许为完成任务
所需的合法身体运动，身体稳定进入目标约束。
UMI-on-Legs 提供操作策略与全身控制分层的参照\cite{umilegs}；
改变 WBC 的系统比较与固定 WBC 的高层方法比较分别报告。

\subsection{Strong Baselines and Data Permissions}
\textbf{受控校准协议 A。} 全部方法获得相同源数据、公开任务目标、
已知身体模型、可用传感与低层控制接口。目标仅提供预算内短探针；不提供
完整任务成功演示或目标在线 RL。校准后的模拟训练只能使用由公开模型和
预算内记录构建的模型。隐藏目标环境仅提供计费探针及封存评估，
实际机制参数不向训练开放。

\textbf{冻结主表。} 主表固定为九行，并要求每行参加四个任务：本文方法、
matched task-only、FADA task-adapted、ACP task-adapted、
Human2Sim2Robot source-trajectory adaptation、DIAL-MPC task-adapted、
ordinary response MPC、public geometry/feedback 与 public source prior。
这一组合覆盖 2025/2026 年近期强方法、Karen Liu 与 Guanya Shi 参与的路线，
并保留三个不依赖论文适配细节的强公共控制\cite{human2sim,fada,acp,dialmpc}。

\begingroup\small
\begin{longtable}{L{3.2cm}L{2.2cm}L{2.7cm}L{4.5cm}}
\toprule
主表行 & 权限 & 身份 & 可解释的比较范围\\\midrule\endfirsthead
\toprule 主表行 & 权限 & 身份 & 可解释的比较范围\\\midrule\endhead
Ours gated transport & B-log & $C+H+GR$ & 与 task-only/FADA 做同权限任务比较；gate 机制另由六格消融识别。\\
Task-only & B-log & $C$ & 同容量、同数据、同训练流与同 checkpoint fallback 的首要对照。\\
FADA & B-log & Planner/IDM task adaptation & 保留公开 planner/IDM 核心；标为任务适配，不冒充原人形定量复现\cite{fada}。\\
ACP & C-supervised & virtual target/stiffness adaptation & 目标示范与力监督单列；不能与 B-log 作相同信息因果比较\cite{acp}。\\
Human2Sim2Robot & C-supervised-model & robot-source trajectory proxy & 原人类资产不可用时明确标为 source-trajectory 改编，不推广成原方法结论\cite{human2sim}。\\
DIAL-\newline MPC & B-log\newline model & annealed diffusion\newline online MPC & 保留在线 candidate rollout 与 receding horizon；模型调用与 deadline miss 进入资源列\cite{dialmpc}。\\
Ordinary response MPC & B-log\newline model & online Adam shooting & 与 DIAL-MPC 使用同一公共 response model 和 decoder/WBC。\\
Public geometry & no-target-learning & 固定几何/反馈 & 检验任务是否已可由强公共反馈解决。\\
Public source prior & no-target-learning & 冻结 source skill & 检验不做目标适配时的直接迁移与部署前拒绝。\\
\bottomrule
\end{longtable}
\endgroup

Play2Perfect、ASAP、SPI-Active 与 XMoP 保留为相关或后续扩展路线，
不在冻结九行中用不完整代理占据主表位置\cite{play2perfect,asap,spi,xmop}。
原生复现、完整系统比较与共同接口改编不能互相替代；每个 task-adapted 行
均报告 source、known-nominal 与目标三段漏斗。

\textbf{部署前拒绝。} 一个实现若耗尽预声明机会并在 source/known-nominal
资格门失败，仍覆盖该任务单元。其每个计划请求记
constrained success $=0$、constraint violation $=1$、capped normalized
time $=1$、pre-deployment rejection $=1$、target execution $=0$。
这里的 violation 是算法部署合同失败，不表示未执行目标轨迹发生了物理超限。
主表同时报告 rejection rate、target execution rate 与物理执行数。

DR/recurrent PPO、经典 Residual RL、通用校准 MPC、Mirage 和 ASID
保留为基础参照\cite{dr,ppo,rppo,residualrl,mjpc,mirage,asid}。
FORGE 是 RA-L 2025 接触参照，不列作已核验的近期顶会主会方法\cite{forge}。
同容量 task-only、deterministic transport、reweighting-only 和去响应/模式/
风险项属于内部消融。强对照可使用连续动作、技能 chunk、历史条件与合理
网络容量，不限制在有限映射候选库。

\textbf{共同权限的完整适应比较。} B-log 向本文方法和 FADA/ASAP 等对照
开放相同总时长、可测信息、reset 与执行约束下的完整自采 action/observation
轨迹，不含目标 reward 或专家示范，各方法自行选择采集策略。
另做共享同一记录缓冲的控制实验，区分采集与适应收益。
ASAP 所需状态对齐与额外测量权限单列；同组方法获得相同可用信息。
不使用 reward 不会使完整目标轨迹自动变成 A 短探针。
B-online 允许计费的目标任务交互与 reward；
C-supervised 比较向本文扩展版本提供相同可用示范并标明监督；信息不同的
原生系统另报性能--成本前沿，不作同预算排名。外部预训练单列。
D-oracle 仅作
具有目标真参数的封存诊断。需要大量随机目标数据的 latent alignment、
目标演示训练的 diffusion/latent-action policy，以及目标在线 residual RL
分别报告数据曲线\cite{lsa,lad}。充分目标训练的策略与真值规划器不并入
少量校准主表，也不被解释为已证明的性能上界。

\textbf{接口与训练公平性。} 受控比较共享传感、可用历史、公开任务 reward、
控制周期、命令边界和安全判据。需要特权 critic 的方法只能读取其训练模型的
状态，隐藏目标参数不得进入训练。原生实现的复现检查与统一任务接口的改编
结果分别记录，不能把本项目同名策略直接视为原论文复现。

\subsection{Budget, Generalization, and Statistical Evaluation}
目标访问按秒数、transitions、reset 和数据类型报告。
短探针子协议可采用 $B\in\{0,5,10,30\}$，每次时长与 fit/validation
分配在正式测试前冻结。需要约两分钟动作/观测记录的 FADA 等方法，
另包含原生时长的数据点\cite{fada}，不能用不合适的极短缓冲代表完整性能。
目标轨迹即使不含 reward，也须声明是否为完整任务执行；原生示范另列。
接口调整、验证与失败尝试全部计费，主动方法受相同总访问预算限制。

训练模拟 transitions、离线更新、调参次数、GPU 时间和部署在线模型调用
另列。比较同模拟量的学习曲线与固定时延下的表现；不同优化器的更新次数
不直接视为相同计算量。MPC 超时须按预先固定的执行规则记入结果。

测试分别留出初态、机制组合与身体。UR5e 的两项协议分开：
冻结 transport 仅拟合 context；以及允许在校准模型内继续更新策略的
few-shot adaptation。已知目标形态上的训练不等于未见身体的冻结泛化。
仿真任务结果与实机结果单列；平台与传感接口未确定时不预设实机结论。

正式评估冻结为八个独立训练/校准实例、两个目标机制组和配对请求。
Task 1--3 每实例为 $10$ 个不同物理初态 $\times2$ 个请求模式
$\times2$ 个机制组，共 $40$ 个请求；Task 4 再交叉两个 body/support cells，
共 $80$ 个请求。确定性公共控制只保留一个方法身份，在同一 request package
上执行，不伪造八个训练重复。所有方法共享相同 task/mechanism/body/initial/mode
键，失败不替换。完整 post-reset 物理与控制器状态 hash 才定义物理初态；
episode 标签或相同 deterministic reset 上的 policy seed 不构成新初态。

八实例由开发数据上的预声明配对 discordance power simulation 冻结。
中等 discordance 场景达到目标 power；高训练异质性仍可能不足，作为限制报告，
不能在查看最终结果后追加实例。推断以训练/校准实例为最高随机层，再对
mechanism/body/initial request block 做 $10{,}000$ 次层级配对 bootstrap。
每个研究问题预声明一个主对比，RQ family 内使用 Holm 修正；报告效应量、
区间和原始分子/分母，不给出无权跨任务总排名。全部支持失败、验证拒绝、
超力、超时和身体失败进入整体分母。

\subsection{Aligned Six-Cell Ablation}
六格消融与主方法使用相同 warm checkpoint、数据流、优化器、更新数、
held-out 模型代理和 task-only $+0.05$ fallback：
\begin{center}
\begin{tabular}{lll}
\toprule
单元 & 目标 & 识别作用\\\midrule
frozen-warm & 不继续训练 & continued training 控制\\
task-only & $C$ & task-first 基准\\
matching-no-band & $C+R$ & response matching 主效应\\
band-only & $C+H$ & additive task band 主效应\\
ungated-full & $C+H+R$ & 无 gate 的 $2\times2$ 完整角\\
gated-full & $C+H+GR$ & selective gate\\
\bottomrule
\end{tabular}
\end{center}
$C,C+R,C+H,C+H+R$ 是严格的无 gate $2\times2$；只用
gated-full$-$ungated-full 识别 gate。六格都参加四任务记账；若整个任务族在
冻结资格门失败，则六格均按部署前拒绝进入分母，不训练或查看目标结果。
历史上用 hinge $H$ 替换 $C$ 的五格工件与当前方法身份不一致，不进入推断。

\subsection{Performance Evidence and Contribution Tests}
主要评价无碰撞且满足力/身体约束的任务成功、失败风险、完成时间，以及
获得这些表现所需的目标数据与计算成本。dense reward 用于训练；统一的任务
成功判据用于评估，累计 reward 不能跨任务代替成功与安全表现。

\begin{enumerate}
\item \textbf{交互效率。} 在同等目标访问权限下，比较成功率--预算曲线，
以及达到预声明成功水平所需的目标交互量。
\item \textbf{接触与恢复。} 在相同约束下比较接触定位、卡滞后的整体恢复、
力峰值与失效概率，并检验未见机制组合。
\item \textbf{任务--行为权衡。} 首先控制任务成本，再比较有任务意义的模式
覆盖、响应和条件对应。无模式请求时允许 baseline 选择更可靠的模式；
有明确模式请求时，向全部方法提供相同条件，并分别评价按请求完成和自主完成。
\item \textbf{摊销执行。} 在与校准模型 MPC 相当的任务表现下，检验 transport
是否降低端到端在线时延；WBC、状态处理与模型计算均计入。
\end{enumerate}
源--目标响应距离下降、额外噪声或更大模型本身不构成性能贡献。
方法是否成立取决于上述比较，不能预设所有任务或指标均优于现有方法。


\section{Discussion}
论文的实证贡献取决于与成熟方法的比较：在相同目标数据权限下改善受约束任务表现、降低达到同等表现所需的目标交互，或在任务表现相当时改善有任务意义的行为选择与执行成本。源响应距离下降本身不足以建立贡献。

若收益仅来自接口增益校准，则应定位为 action adaptation；若 task-aware ASID 已达到相同效果，则需要谨慎评估 transport 的增量价值；若 ordinary response prediction 与完整方法等效，则缩小 causal 主张。目标分布不必全局复制源分布，必要的模式改变应由可行性和任务结果解释。

\section{Limitations}
\begin{enumerate}
\item 校准与任务共享的结构化机制假设、共同状态充分性及实际访问支持决定泛化范围；未知迟滞/接触分支可能破坏模型。
\item task geometry 与 source anchoring 可能引入偏置；不能把 transport objective 最优等同于任意目标 task-only 最优。
\item 短时模型的组合误差、rare mode 概率与 model-based CVaR 需要独立验证；模型风险约束不是实机安全证书。
\item source-conditioned map/flow 不保证保留所有有效 modes，也不保证创造任意新模式；低支持区域需要新增交互或更宽 prior。
\item intervention 支持规定协议下的响应辨识，不恢复唯一 latent DAG 或任意跨身体个体反事实。
\item 首版仿真采用结构化接触状态，实物平台待确定；不能直接外推到通用视觉、任意软体机器人或真实双臂迁移。
\item Task 4 依赖具备声明末端/力接口的 WBC 和共同源支持；同工位局部操作成功不等于学得导航、步态或任意身体的全身操作策略。身体专属可行性与模型误差需单独验收。
\end{enumerate}

\section{Conclusion}
本文希望通过有限干预校准目标执行机制，将源任务条件行为分布转化为可执行目标行为分布。完整首版应包含 source prior、明确的 action decoding 接口、校准响应、摊销 transport、分布与闭环指标，以及同预算已有方法比较。方法与实验决定贡献，理论解释其依据与边界。

\appendix
\section{Supporting Derivations}
\subsection{Identifiability and Fixed-Design Uncertainty}
若两个参数产生相同校准结果，则 $v=\theta_1-\theta_2\in\ker A_{\mathcal B}$。相关响应一致当且仅当 $B_gv=0$，直接得到主文核空间条件。

固定设计且 $\xi\sim\mathcal N(0,\sigma^2I)$ 时
\begin{equation}
 \hat b_g=B_gA_{\mathcal B}^\dagger y,\qquad
 \Sigma_g=\sigma^2B_g(A_{\mathcal B}^\top A_{\mathcal B})^\dagger B_g^\top.
\end{equation}
在可辨识条件下，以概率至少 $1-\eta$，
\begin{equation}
 \|\hat b_g-b_g\|\le
 \sqrt{\tr\Sigma_g}+\sqrt{2\|\Sigma_g\|_{\rm op}\log(1/\eta)}.
\end{equation}
多个查询需联合控制；真实响应模型还需结构偏差、模式概率与噪声 law 的误差。小均值误差不自动推出整个 kernel 接近。

\subsection{Propagation and Objective Error}
对保持同一源参考的真实和预测增广状态 law 应用 triangle inequality：
\begin{equation}
 d_{k+1}\le
 W_1(\mu_k\mathbb K_t^\varphi,\mu_k\widehat{\mathbb K}_t^\varphi)
 +W_1(\mu_k\widehat{\mathbb K}_t^\varphi,\hat\mu_k\widehat{\mathbb K}_t^\varphi)
 \le\varepsilon_{\mathcal B}+Ld_k.
\end{equation}
对 trajectory cost 必须使用逐步 coupling，而非仅边缘 law；在对应的 kernel coupling 和统一 Lipschitz 条件下，
\begin{equation}
 \varepsilon_F\le
 (L_C+\lambda_sL_c)\varepsilon_{\mathcal B}
 \sum_{k=1}^{H}\sum_{j=0}^{k-1}L^j.
\end{equation}
同一个 source coupling、相同 transport 输入及真实/模型共同可行集合为前提。风险 tightening 不成立时，模型可行解未必真实可行。UOT 增强与不连续事件需单独误差分析。

主文目标界通过加减模型目标获得：
\begin{align}
 F_t(\hat\varphi)-F_t(\varphi^*)
 &= [F_t(\hat\varphi)-\widehat F_t(\hat\varphi)]
 +[\widehat F_t(\hat\varphi)-\widehat F_t(\varphi^*)]
 +[\widehat F_t(\varphi^*)-F_t(\varphi^*)]\\
 &\le2\varepsilon_F+\varepsilon_{\rm opt}.
\end{align}
这里优化的是约定 transport 目标，不移除 source regularization 后重新解释为无约束 RL regret。

\begin{thebibliography}{99}
\small
\bibitem{lsa} T. Wang, D. Bhatt, X. Wang, N. Atanasov. Cross-Embodiment Robot Manipulation Skill Transfer using Latent Space Alignment. arXiv:2406.01968, 2024. \url{https://arxiv.org/abs/2406.01968}
\bibitem{mirage} L. Y. Chen et al. Mirage: Cross-Embodiment Zero-Shot Policy Transfer with Cross-Painting. RSS, 2024. \url{https://robot-mirage.github.io/}
\bibitem{fada} A. Xie et al. FADA: Few-Shot Domain Adaptation via Dynamics Alignment for Humanoid Control. CoRL, 2026, accepted; arXiv:2606.28476. \url{https://arxiv.org/abs/2606.28476} 录用与代码：\url{https://github.com/LeCAR-Lab/FADA}
\bibitem{lad} E. Bauer, E. Nava, R. K. Katzschmann. Latent Action Diffusion for Cross-Embodiment Manipulation. ICRA, 2026; arXiv:2506.14608v4. \url{https://arxiv.org/abs/2506.14608}
\bibitem{xskill} M. Xu et al. XSkill: Cross Embodiment Skill Discovery. CoRL, 2023. \url{https://xskill.cs.columbia.edu/}
\bibitem{human2sim} T. G. W. Lum, O. Y. Lee, K. Liu, J. Bohg. Crossing the Human-Robot Embodiment Gap with Sim-to-Real RL using One Human Demonstration. CoRL, 2025. \url{https://proceedings.mlr.press/v305/lum25a.html}
\bibitem{dr} X. B. Peng et al. Sim-to-Real Transfer of Robotic Control with Dynamics Randomization. ICRA, 2018. \url{https://arxiv.org/abs/1710.06537}
\bibitem{gat} J. P. Hanna, P. Stone. Grounded Action Transformation for Robot Learning in Simulation. AAAI, 2017. \url{https://ojs.aaai.org/index.php/AAAI/article/view/11044}
\bibitem{rma} A. Kumar et al. RMA: Rapid Motor Adaptation for Legged Robots. RSS, 2021. \url{https://arxiv.org/abs/2107.04034}
\bibitem{asid} M. Memmel et al. ASID: Active Exploration for System Identification in Robotic Manipulation. ICLR, 2024. \url{https://weirdlabuw.github.io/asid/}
\bibitem{spi} N. Sobanbabu et al. Sampling-based System Identification with Active Exploration for Legged Sim2Real Learning. CoRL, 2025. \url{https://proceedings.mlr.press/v305/sobanbabu25a.html}
\bibitem{optimalexploration} A. Wagenmaker, G. Shi, K. G. Jamieson. Optimal Exploration for Model-Based RL in Nonlinear Systems. NeurIPS, 2023. \url{https://proceedings.neurips.cc/paper_files/paper/2023/hash/31e018f43ab9c7065c058cc2c5848128-Abstract-Conference.html}
\bibitem{otreal} S. Cheng, L. Ma, Z. Chen, A. Mandlekar, C. Garrett, D. Xu. Generalizable Domain Adaptation for Sim-and-Real Policy Co-Training. NeurIPS, 2025. \url{https://arxiv.org/abs/2509.18631}
\bibitem{densityflow} X. Duan, D. Chen. Trajectory-Optimized Density Control with Flow Matching. arXiv:2510.06666, 2025. \url{https://arxiv.org/abs/2510.06666}
\bibitem{dara} J. Liu, H. Zhang, D. Wang. DARA: Dynamics-Aware Reward Augmentation in Offline Reinforcement Learning. ICLR, 2022. \url{https://arxiv.org/abs/2203.06662}
\bibitem{h2oplus} H. Niu et al. H2O+: An Improved Framework for Hybrid Offline-and-Online RL with Dynamics Gaps. ICRA, 2025. \url{https://arxiv.org/abs/2309.12716}
\bibitem{transic} Y. Jiang et al. TRANSIC: Sim-to-Real Policy Transfer by Learning from Online Correction. CoRL, 2024. \url{https://arxiv.org/abs/2405.10315}
\bibitem{isaaclabtasks} Isaac Lab contributors. Available Environments: Factory, FORGE, and AutoMate. Official documentation, accessed 2026-10-12. \url{https://isaac-sim.github.io/IsaacLab/main/source/overview/environments.html}
\bibitem{umilegs} H. Ha, Y. Gao, Z. Fu, J. Tan, S. Song. UMI-on-Legs: Making Manipulation Policies Mobile with Manipulation-Centric Whole-body Controllers. CoRL, 2024; PMLR 270:5254--5270, 2025. \url{https://proceedings.mlr.press/v270/ha25a.html}
\bibitem{menagerie} MuJoCo Menagerie contributors. MuJoCo Menagerie: Robot model collection. Repository, accessed 2026-10-12. \url{https://github.com/google-deepmind/mujoco_menagerie}
\bibitem{ppo} J. Schulman et al. Proximal Policy Optimization Algorithms. arXiv:1707.06347, 2017. \url{https://arxiv.org/abs/1707.06347}
\bibitem{rppo} SB3-Contrib contributors. Recurrent PPO: Official implementation documentation. Accessed 2026-10-14. \url{https://sb3-contrib.readthedocs.io/en/master/modules/ppo_recurrent.html}
\bibitem{forge} M. Noseworthy et al. FORGE: Force-Guided Exploration for Robust Contact-Rich Manipulation under Uncertainty. IEEE Robotics and Automation Letters, 2025; arXiv:2408.04587. \url{https://arxiv.org/abs/2408.04587}
\bibitem{mjpc} T. Howell et al. Predictive Sampling: Real-time Behaviour Synthesis with MuJoCo. arXiv:2212.00541, 2022. \url{https://arxiv.org/abs/2212.00541}
\bibitem{residualrl} T. Johannink et al. Residual Reinforcement Learning for Robot Control. ICRA, 2019. \url{https://arxiv.org/abs/1812.03201}
\bibitem{play2perfect} T. G. W. Lum, K. Kedia, C. K. Liu, J. Bohg. Play2Perfect: What Matters in Dexterous Play Pretraining for Precise Assembly? CoRL, 2026, accepted; arXiv:2606.26428. \url{https://arxiv.org/abs/2606.26428} 录用公告：\url{https://iprl.stanford.edu/}
\bibitem{asap} T. He et al. ASAP: Aligning Simulation and Real-World Physics for Learning Agile Humanoid Whole-Body Skills. RSS, 2025. \url{https://roboticsproceedings.org/rss21/p066.html}
\bibitem{acp} Y. Hou et al. Adaptive Compliance Policy: Learning Approximate Compliance for Diffusion Guided Control. ICRA, 2025. \url{https://adaptive-compliance.github.io/}
\bibitem{dialmpc} H. Xue, C. Pan, Z. Yi, G. Qu, G. Shi. Full-Order Sampling-Based MPC for Torque-Level Locomotion Control via Diffusion-Style Annealing. ICRA, 2025. \url{https://ieeexplore.ieee.org/document/11127320/}
\bibitem{xmop} P. K. Rath, N. Gopalan. XMoP: Whole-Body Control Policy for Zero-Shot Cross-Embodiment Neural Motion Planning. ICRA, 2025. \url{https://prabinrath.github.io/xmop/}
\end{thebibliography}
\end{document}
